% Generated draft: requires mathematical and source review.
\documentclass[11pt]{article}
\usepackage[margin=25mm]{geometry}
\usepackage{fontspec}
\setmainfont{Latin Modern Roman}
\usepackage{amsmath,amssymb,mathtools,unicode-math}
\setmathfont{Latin Modern Math}
\usepackage{enumitem,needspace,xurl,hyperref}
\setlength{\parindent}{0pt}
\setlength{\parskip}{5pt}
\setlength{\emergencystretch}{3em}
\newcommand{\N}{\mathbb{N}}
\newcommand{\Z}{\mathbb{Z}}
\newcommand{\Q}{\mathbb{Q}}
\newcommand{\R}{\mathbb{R}}
\newcommand{\C}{\mathbb{C}}
\newcommand{\F}{\mathbb{F}}
\newcommand{\A}{\mathbb{A}}
\newcommand{\PP}{\mathbb{P}}
\newcommand{\E}{\mathbb{E}}
\newcommand{\eps}{\varepsilon}
\newcommand{\dd}{\,\mathrm{d}}
\begin{document}
\section*{UnsolvedMath ID: 10600027}
\textbf{Source problem code:} catalogue entry
\subsection{Definitions and mathematical statement}
Let \(\mathcal V\) be virtual knots modulo the intended virtual Reidemeister and detour moves. In the usual oriented movie description, virtual cobordisms also allow births, deaths and saddle moves, and concordance is represented by a connected genus-zero cobordism between the two knots. The exact allowed surface/stabilization and framing conventions must be fixed.

The source asks for concordance and cobordism invariants of virtual knots, and specifically a classification of
\[
 \mathcal V/\!\sim_{\mathrm{pass}},
\]
where \(\sim_{\mathrm{pass}}\) is generated by virtual equivalence and the direct virtual generalization of the classical pass move intended by the source. A complete classification should give explicit realizable invariant values that are equal precisely on pass-equivalent knots.

The motivating classical comparison is pass equivalence classified by the Arf invariant. For a classical knot, this is the \(\Z/2\Z\)-valued Arf invariant of the quadratic form obtained from a Seifert form modulo \(2\). The problem does not assert that the virtual quotient has only two classes or that this classical construction automatically extends to all virtual knots. No pass-move diagram is supplied, so the precise virtual relation cannot be recovered solely from its name.
\subsection{Short English statement}
Develop virtual-knot concordance and cobordism invariants, especially a complete classification under the intended virtual pass moves, using the classical Arf classification as motivation.
\subsection{Sources}
\begin{itemize}
\item \emph{Unsolved Problems in Virtual Knot Theory and Combinatorial Knot Theory}, revised survey, problem-list item 27. \url{https://arxiv.org/abs/1409.2823}. Recorded access 2026-07-29; extraction source-tex. No individual proposer, pass-move figure or separate reference is supplied.
\item Local UnsolvedMath registry, \texttt{data/unsolvedmath/problems.json}, permanent ID \(10600027\), exact source code \texttt{catalogue entry}. Its partially\_solved/PARTIAL-PROGRESS summary mentions 2016 and 2024 studies of pass/Arf-type invariants and active concordance work, without complete citations or a documented full classification.
\end{itemize}
\end{document}
